\section{Introdução}

O monitoramento da qualidade da água em áreas remotas da Amazônia é limitado pela infraestrutura esparsa, pelo alto custo logístico das campanhas de campo e pela conectividade de rede intermitente. A implantação de Edge AI é a motivação central deste trabalho: a inferência local, totalmente offline, é necessária para manter o monitoramento funcional sob comunicações instáveis e orçamentos de energia limitados. Nesse contexto, a dependência contínua de inferência em nuvem é operacionalmente frágil e energeticamente cara, especialmente em implantações alimentadas por bateria e campanhas de longa duração. Ao mesmo tempo, variáveis de qualidade da água como pH, oxigênio dissolvido (OD) e turbidez são altamente relevantes para a triagem ambiental e a mitigação de riscos à saúde pública.

Pipelines convencionais de aprendizado de máquina podem oferecer forte desempenho preditivo, mas muitos são projetados para execução centrada em servidores e não tratam diretamente as restrições energéticas na borda. Modelos neuromórficos baseados em redes neurais pulsadas (SNNs, do inglês \textit{spiking neural networks}) oferecem uma alternativa orientada a eventos, compatível com sistemas embarcados de baixo consumo; contudo, evidências práticas em fluxos reais de sensoriamento ambiental ainda são limitadas, particularmente quando se combinam predição multimodal e avaliação energética em nível de plataforma.

Este artigo propõe e avalia um pipeline SNN multimodal offline para predição da qualidade da água em um cenário de campo amazônico, incluindo implantação de ponta a ponta e análise energética multiplataforma. O estudo integra processamento de dados, treinamento de modelos, interpretação geoespacial e inferência embarcada em hardware heterogêneo (ESP32, Raspberry Pi e dispositivos Android).

As principais contribuições deste trabalho são \rev{quatro}: (i) o desenvolvimento de um pipeline de predição da qualidade da água baseado em SNN multimodal, projetado para implantação offline na borda; (ii) a criação de um conjunto de dados georreferenciado aumentado, derivado de medições de campo no Lago do Serpa; (iii) uma avaliação comparativa de latência--energia em plataformas heterogêneas, incluindo ESP32, Raspberry Pi e dispositivos Android\rev{; e (iv) um protocolo de avaliação estatisticamente fundamentado para regressão ambiental com poucos dados, combinando avaliação com sementes repetidas, intervalos de confiança e testes de significância não paramétricos. Para além da avaliação energética, a contribuição científica reside em demonstrar, de ponta a ponta, que a regressão neuromórfica multimodal de alvos físico-químicos de qualidade da água é viável desde a aquisição em campo até a inferência totalmente offline em três classes distintas de hardware de borda --- uma combinação ainda não reportada para o monitoramento de água doce na região amazônica.}
